\documentclass[11pt]{article}

% Readable single-column lesson and solutions.
\usepackage[letterpaper,margin=0.4in]{geometry}
\usepackage[T1]{fontenc}
\usepackage{microtype}
\usepackage{amsmath,amssymb}
\usepackage{enumitem}
\usepackage{needspace}
\usepackage{booktabs}

\pagestyle{plain}
\linespread{1.08}
\setlength{\parindent}{0pt}
\setlength{\parskip}{4pt}
\setlength{\emergencystretch}{2em}
\setlist[enumerate]{leftmargin=*,itemsep=5pt,topsep=4pt,parsep=0pt}
\newcommand{\stageheading}[1]{%
  \par\Needspace{6\baselineskip}%
  \addvspace{.65\baselineskip}%
  {\centering\large\bfseries #1\par}%
  \nopagebreak[4]}
\newcommand{\solution}[1]{%
  \par\Needspace{5\baselineskip}\addvspace{.8\baselineskip}%
  {\large\bfseries Solution #1\par}\nopagebreak[4]}
\clubpenalty=10000
\widowpenalty=10000
\displaywidowpenalty=10000

\begin{document}

% Compact opening only; restore the original geometry and typography for solutions.
\newgeometry{left=0.65in,right=0.65in,top=0.25in,bottom=0.29in}
\begingroup
\linespread{1}\selectfont
\setlength{\parskip}{2pt}
\setlength{\abovedisplayskip}{5pt plus 1pt minus 1pt}
\setlength{\belowdisplayskip}{5pt plus 1pt minus 1pt}
\setlength{\abovedisplayshortskip}{3pt plus 1pt}
\setlength{\belowdisplayshortskip}{3pt plus 1pt}
\setlist[enumerate]{leftmargin=*,itemsep=3pt,topsep=3pt,parsep=0pt}
\renewcommand{\stageheading}[1]{%
  \par\Needspace{3\baselineskip}%
  \addvspace{.35\baselineskip}%
  {\centering\large\bfseries #1\par}%
  \nopagebreak[4]}

\vspace*{0pt}
{\centering\Large\bfseries Permutations and Combinations I\par}

\stageheading{Introduction: Counting Choices}

\textbf{Multiplication principle.} Count a process one step at a time.
If the number of available choices at each step is the same for every
possible sequence of earlier choices, multiply these numbers to count
the complete outcomes. \emph{Order matters} means that exchanging
positions changes the outcome.

\textbf{Notation and when to use it.} Let $n$ be the total number of
objects, a nonnegative integer, and let $r$ be the number selected,
an integer with $0\leq r\leq n$.
\begin{itemize}[leftmargin=*,itemsep=3pt,topsep=3pt,parsep=0pt]
\item To arrange all $n$ distinct objects in a row, there are $n$ choices
for the first position, then $n-1$, and so on. Their product is called
\emph{factorial}: $n!=n(n-1)\cdots2\cdot1$ for $n\geq1$.
By convention, $0!=1$: there is one empty arrangement.
\item To fill $r$ ordered positions with distinct objects chosen from $n$
without replacement, the choices decrease from $n$ to $n-r+1$.
The count is $\dfrac{n!}{(n-r)!}$; each object can be used at most once.
\item To choose $r$ distinct objects when order does not matter, divide
the ordered count by the $r!$ orders of each selected group. The count is
$\displaystyle\binom{n}{r}=\dfrac{n!}{r!(n-r)!}$, read ``$n$ choose $r$.''
\item To arrange all $n$ objects when some copies are indistinguishable,
first imagine labelling every copy. Then divide by the orders of labels
within each type, since these exchanges leave the arrangement unchanged.
For $k$ types with $m_i$ copies of type $i$, the count is
$\dfrac{n!}{m_1!m_2!\cdots m_k!}$.
Here $k$ is a positive integer, $i=1,\ldots,k$ labels the types, each
$m_i$ is a positive integer, and $m_1+\cdots+m_k=n$.
\end{itemize}

\stageheading{Worked Examples}

\textbf{Worked example 1: Three distinct roles.}
Five people volunteer for an event. Assign a coordinator, a treasurer,
and a secretary, with no person holding more than one role. Order matters:
exchanging the coordinator and treasurer changes who does each job.
Choose the coordinator in five ways, the treasurer in four remaining
ways, and the secretary in three remaining ways. Multiplying gives
\[
5\cdot4\cdot3=\frac{5!}{(5-3)!}=\boxed{60}.
\]

\textbf{Worked example 2: A committee without roles.}
Choose three of the same five volunteers for a committee with no assigned
roles. Order does not matter: exchanging the positions in a list of
members leaves the committee unchanged. Each committee corresponds to
$3!=3\cdot2\cdot1=6$ ordered role assignments: its three members can
fill the coordinator, treasurer, and secretary roles in six ways.
Thus the previous count includes each committee six times. Dividing gives
\[
\binom{5}{3}=\frac{5!}{3!(5-3)!}=\frac{60}{3!}
=\frac{60}{6}=\boxed{10}.
\]


\Needspace{3\baselineskip}
\textbf{Worked example 3: Together, but with a forbidden pair.}
Seven distinct runners, $P,Q,R,S,T,U,V$, line up for a photograph.
How many rows put $P$ next to $Q$, but keep $R$ away from $S$?
Either order within a pair is allowed. All seven runners appear once;
order matters because the photograph has a left and a right.

\textbf{First enforce the required adjacency.} Replace $P,Q$ by a
single block. We now arrange six distinct units: the block and the other
five runners. The block has two internal orders, $PQ$ and $QP$.
For example, contraction and expansion relate these descriptions:
\[
T\quad\boxed{QP}\quad R\quad V\quad S\quad U
\quad\longleftrightarrow\quad
T\quad Q\quad P\quad R\quad V\quad S\quad U.
\]
There are $6!$ orders of the units and two choices inside the block,
so $2\cdot6!=1440$ rows meet the required adjacency.

\textbf{Then subtract the forbidden adjacency.} Among those rows,
the unwanted ones also put $R,S$ together. Contract both pairs.
\begin{center}
\begin{tabular}{@{}lll@{}}
\toprule
Rows counted & Units to arrange & Number \\
\midrule
$P,Q$ together & $[PQ],R,S,T,U,V$ & $2\cdot6!=1440$ \\
Both pairs together & $[PQ],[RS],T,U,V$ & $2\cdot2\cdot5!=480$ \\
\bottomrule
\end{tabular}
\end{center}
The brackets denote pairs in either internal order. Thus the answer is
\[
2\cdot6!-4\cdot5!=1440-480=\boxed{960}.
\]
Every row with $P,Q$ adjacent has a unique contracted row and a unique
internal order. The same is true when both pairs are contracted, even
if their blocks touch. Expanding recovers the original row, so each
count includes each relevant row exactly once. The unwanted rows form
a subset of the first count; subtracting leaves exactly the valid rows.

\stageheading{1. Progressive Problems}

\begin{enumerate}[label=\textbf{\arabic*.}]
\item Six books have different titles. In how many ways can all six be
placed in a single row on a shelf? Each book is used once, and two rows
are different if their left-to-right orders differ.

\item A club has nine distinct members. In how many ways can it choose
a committee of four different members? There are no assigned roles, and
the order in which members are chosen does not matter.

\item Eight distinct students enter a contest with three awards: first,
second, and third place. In how many ways can the awards be assigned if
each award goes to one student and no student receives more than one award?

\item In how many different left-to-right strings can all the letters of
\texttt{BALLOON} be arranged? Use each letter exactly as often as it appears
in the word. Copies of the same letter are indistinguishable, and the
strings need not be English words.

\item Arrange all the letters of \texttt{BALLOON} in a left-to-right
string, but do not allow the two L's to be adjacent. Copies of the same
letter are indistinguishable; each letter must be used as often as it
appears in the word. How many strings are possible? Explain how to place
the repeated letters without counting a string more than once.

\item Six distinct books, labelled $A,B,C,D,E,F$, are placed in a row,
each exactly once. How many left-to-right arrangements have $A$ and $B$
next to each other, but $C$ and $D$ not next to each other? Either order
within an adjacent pair is allowed.

\item Four girls, Ada, Bella, Cleo, and Dina, and three boys, Eli, Finn,
and Gus, are available for a photograph. Choose five different students
and arrange them in a row. The row must include Ada, at least two girls,
and at least two boys, and no two boys may stand next to each other.
How many different left-to-right rows are possible? All seven students
are distinguishable, and each selected student appears exactly once.
\item Three different numbers are chosen from $\{1,2,\ldots,10\}$.
No two chosen numbers may differ by $1$. How many unordered sets can be
chosen? Each number can be used at most once.

For a useful representation, write a chosen set as $a<b<c$.
Find a reversible transformation that removes the compulsory spaces
between these numbers, and use it to count the sets.
\end{enumerate}

\stageheading{2. Olympiad Challenges}
\begin{enumerate}[label=\textbf{\arabic*.},start=9]
\item Choose an unordered five-element subset of $\{1,2,\ldots,12\}$
with exactly two consecutive pairs. Here a consecutive pair means
$\{i,i+1\}$ with both numbers selected. Overlap is allowed: selecting
$4,5,6$ creates two pairs, $\{4,5\}$ and $\{5,6\}$.
How many such subsets are there? Explain how to include overlapping
and separate pairs without double counting.

\item Arrange the numbers $1,2,\ldots,7$ in a row, each exactly once.
Call an index $i\in\{1,2,\ldots,6\}$ a \emph{success} if $i$ is
immediately followed by $i+1$, in that order. How many permutations
have exactly two successes? Successes may overlap: for instance,
the consecutive segment $3,4,5$ contributes successes at both $i=3$
and $i=4$. Two permutations are different if their left-to-right orders differ.

\item A robot travels from $(0,0)$ to $(6,6)$ using exactly six steps
$R=(1,0)$ and six steps $U=(0,1)$, in some order. At every point of
its route it must satisfy $y\leq x$: it may touch the diagonal $y=x$
but may never go above it. How many different routes are possible?
Two routes are different if their step sequences differ. Find a
reversible way to count the forbidden routes rather than listing all
allowed routes.
\end{enumerate}

\clearpage
\endgroup
\restoregeometry
\stageheading{3. Solutions}

\solution{1}
Order matters because the shelf has a left-to-right order. There are six
choices for the first position, five for the second, and so on. Every
book is distinct and used once, so there is no division for identical
objects and no repetition of choices. The multiplication principle gives
 \vspace{-3mm}
\[
6\cdot5\cdot4\cdot3\cdot2\cdot1=6!=\boxed{720}.
\]
\vspace{-12mm}



\solution{2}
A committee is an unordered group. First count ordered selections of
four different members: there are $9\cdot8\cdot7\cdot6$ choices.
Every committee appears $4!$ times in this count, once for each order of
its four members. Dividing by this common overcount gives
\[
\binom94=\frac{9\cdot8\cdot7\cdot6}{4!}=\boxed{126}.
\]

\vspace{-8mm}

\solution{3}
The awards are distinct positions, so order matters: exchanging the first-
and second-place winners changes the result. Choose the first-place
winner in eight ways, the second in seven, and the third in six.
The decreasing numbers of choices enforce the rule that a student cannot
win twice. There is no division by $3!$, since the three award assignments are distinct. Thus
\[
8\cdot7\cdot6=\frac{8!}{5!}=\boxed{336}.
\]
 \vspace{-8mm}



\solution{4}
There are seven letter positions. Temporarily distinguish the two L's
as $L_1,L_2$ and the two O's as $O_1,O_2$. These seven labelled objects
have $7!$ ordered arrangements, with no object used twice.

In an actual string, exchanging the labels on the L's or on the O's
makes no visible change. Every actual string therefore corresponds to
exactly $2!\,2!$ labelled arrangements. All other letters occur once. Order still matters; only exchanges of indistinguishable copies are removed so
\[
\frac{7!}{2!\,2!}=\boxed{1260}.
\]
 \vspace{-12mm}

\solution{5}
We count ordered strings, but the two L's are indistinguishable, as are
the two O's. Build the string by placing the letters other than L first.
The five letters $B,A,O,O,N$ have $5!/2!=60$ distinct orders: labelling
the O's temporarily counts each visible order exactly twice.

Every such order has six gaps, including its two ends. For example,
\[
\underbrace{\square}_{1}\ B\
\underbrace{\square}_{2}\ A\
\underbrace{\square}_{3}\ O\
\underbrace{\square}_{4}\ O\
\underbrace{\square}_{5}\ N\
\underbrace{\square}_{6}.
\]
Put the two L's in two different gaps. This is necessary and sufficient
for separation: the letters of the base row separate different gaps,
whereas two L's in the same gap would touch. Choose the occupied gaps in
$\binom62=15$ ways. Do not multiply by $2!$, because exchanging the L's
makes no new string. The answer is
\[
\frac{5!}{2!}\binom62=60\cdot15=\boxed{900}.
\]
Removing the L's from a valid string uniquely recovers its base order
and its two occupied gaps. Inserting them reverses this procedure, so
nothing is missed or counted twice, even when the O's are next to one
another.

\textbf{Independent check by subtraction.} There are $1260$ unrestricted
strings from Solution 4. If the L's touch, replace them by the single
unit $[LL]$. The six units $[LL],B,A,O,O,N$ have $6!/2!=360$ orders.
The block has only one internal order because its letters are identical.
Thus $1260-360=900$, agreeing with the gap count.

\solution{6}
Count all rows with $A,B$ together, then subtract those in which $C,D$
are also together. Treat $A,B$ as one block. Together with $C,D,E,F$,
this gives five distinct objects, arranged in $5!$ ways. The block can
be $AB$ or $BA$, giving $2\cdot5!$ rows.

Each row has a unique decomposition into these blocks, even when the
blocks touch, so there is no overcount. All six books remain distinct
and appear once; both their order and each block's internal order matter.

For the unwanted rows, treat both pairs as blocks. The four objects
are the two blocks, $E$, and $F$. Each block has two internal orders,
so there are $2\cdot2\cdot4!$ unwanted rows. Hence the required count is
\[
2\cdot5!-4\cdot4!=240-96=\boxed{144}.
\]

 \vspace{-10mm}
\solution{7}
Since five students are selected with at least two of each gender, there
are exactly two possible compositions: two girls and three boys, or
three girls and two boys. These cases are disjoint and exhaust the possibilities.
In each case, choose the students first, then arrange them; no student
is chosen twice.

The table separates selection (unordered) from placement (ordered).
Ada is already included in each girls' selection.
\begin{center}
\renewcommand{\arraystretch}{1.2}
\begin{tabular}{@{}ccccc@{}}
\toprule
Girls, boys & Choose girls & Choose boys & Arrange using gaps & Total \\
\midrule
$2,3$ & $\binom31$ & $\binom33$ & $2!\binom33 3!$ & $36$ \\
$3,2$ & $\binom32$ & $\binom32$ & $3!\binom42 2!$ & $648$ \\
\bottomrule
\end{tabular}
\end{center}
Each square can hold at most one boy. The symbols $G_1,G_2,G_3$ denote
positions in the chosen girls' order, not additional students.
For a fixed girls' order, the gaps have the following forms:
\[
\begin{array}{ll}
\text{Two girls:}&\square\ G_1\ \square\ G_2\ \square,\\[4pt]
\text{Three girls:}&\square\ G_1\ \square\ G_2\ \square\ G_3\ \square.
\end{array}
\]


\textbf{Case 1: Two girls and three boys.}
Ada must be one of the girls, so choose her companion in $3$ ways.
All three boys must be chosen. Arrange the two girls in $2!$ ways.
They determine three gaps: before, between, and after them. To separate
three boys, put exactly one in every gap, in $3!$ orders. This gives
\[
3\cdot2!\cdot3!=36.
\]

\textbf{Case 2: Three girls and two boys.}
Choose two of the other three girls to join Ada, and choose two of the
three boys. Arrange the girls in $3!$ ways. Their row has four gaps;
choose two different gaps for the boys and assign the boys to them
in $2!$ orders. This gives
\[
\binom32\binom32\,3!\binom42\,2!=648.
\]
Using different gaps is exactly what prevents adjacent boys. Conversely,
removing the boys from any valid row uniquely recovers the girls' order,
the occupied gaps, and the boys' order, so every valid row is counted once.
Selection is unordered, but the final positions are ordered and all
students are distinguishable. Adding the cases gives
\[
36+648=\boxed{684}.
\]



\clearpage
\solution{8}
\textbf{Key insight: remove the compulsory spaces between chosen numbers.}
Write the three chosen numbers in increasing order as $a<b<c$.
This is a unique way to represent an unordered set; it does not add
an ordering choice. The restriction says $b\geq a+2$ and $c\geq b+2$.

Define
\[
x=a,\qquad y=b-1,\qquad z=c-2.
\]
Then $1\leq x<y<z\leq8$. Thus every allowed set produces an ordinary
three-element subset of $\{1,2,\ldots,8\}$.

This correspondence is reversible. From any $1\leq x<y<z\leq8$,
set $a=x$, $b=y+1$, and $c=z+2$. These numbers lie between $1$ and $10$,
and their successive differences are at least $2$. The two constructions
undo each other, so this is a bijection: nothing is lost or counted twice.
For example, the compulsory spacings are removed as follows:
\[
\begin{array}{c|ccc}
 &\text{first}&\text{second}&\text{third}\\\hline
\text{Original choice}&2&5&9\\
\text{Subtract}&0&1&2\\\hline
\text{Compressed choice}&2&4&7
\end{array}
\qquad
(2,5,9)\ \underset{+(0,1,2)}{\overset{-(0,1,2)}{\rightleftarrows}}\ (2,4,7).
\]
The subtraction removes one compulsory unit of separation before the
second entry and two before the third. Actual gaps may be larger:
the transformation does not assume that the chosen numbers differ
by exactly two. All numbers are distinct and none is repeated. Therefore the answer is
\[
\binom83=\frac{8\cdot7\cdot6}{3!}=\boxed{56}.
\]

\vspace{-8mm}
\solution{9}
Order does not matter: we count subsets, each represented by its
increasing list. Split a chosen list into \emph{maximal runs} of
consecutive numbers. For instance, $\{2,3,4,8,11\}$ has runs
$(2,3,4),(8),(11)$. Maximal means that two different runs have at least
one unselected number between them.

\textbf{Translate pairs into runs.} A run of length $\ell$ contributes
$\ell-1$ consecutive pairs, including overlaps. If five selected numbers
form $r$ runs, the total is $5-r$. Exactly two pairs therefore means
exactly three runs. The possible lengths, read from left to right, are:
\begin{center}
\begin{tabular}{@{}lll@{}}
\toprule
Type & Ordered run lengths & Number of orders \\
\midrule
Two overlapping pairs & $(3,1,1),(1,3,1),(1,1,3)$ & $3$ \\
Two separate pairs & $(2,2,1),(2,1,2),(1,2,2)$ & $3$ \\
\bottomrule
\end{tabular}
\end{center}
These are all positive triples summing to five: after giving each run
one number, the two remaining numbers either extend one run twice or
two runs once each. These types are disjoint, and together exhaustive.

\textbf{Place a fixed run-length order.} Of the twelve numbers, five
are selected and seven are not. At least one unselected number must
separate each of the two neighboring runs. Reserve these two separators.
The other five unselected numbers can go into four gaps: \vspace{-3mm}
\[
\underbrace{\circ\cdots\circ}_{g_0}\ [\text{run 1}]\
\underbrace{\circ\cdots\circ}_{1+g_1}\ [\text{run 2}]\
\underbrace{\circ\cdots\circ}_{1+g_2}\ [\text{run 3}]\
\underbrace{\circ\cdots\circ}_{g_3}.
\]
Here $g_0,g_1,g_2,g_3\geq0$ and $g_0+g_1+g_2+g_3=5$.
To count these gap choices, write five identical dots and three dividers;
the dividers split the dots into four groups. Consecutive dividers and
end dividers allow empty groups. Choosing the three divider positions
among eight positions gives $\binom83=56$ choices.

Starting at 1, a run-length order and these gaps uniquely specify which
of the twelve numbers are selected. Conversely, every valid subset has
a unique sequence of maximal run lengths and gaps; subtracting one from
each internal gap recovers the $g_i$. This reversible construction proves
that all six length orders have the same count and do not overlap.
There is no extra factorial for the order of the selected numbers.
Thus
\[
(3+3)\binom83=6\cdot56=\boxed{336}.
\]

\clearpage
\solution{10}
\textbf{Key insight: required consecutive links can be contracted into blocks,
even when they overlap.}
For each $i\in\{1,\ldots,6\}$, let $E_i$ mean that $i$ is immediately
followed by $i+1$. If a specified set of $k$ links is required, merge
each chain of those links into one block. For example, requiring $E_3$
and $E_4$ creates the single block $(3,4,5)$.

For comparison, requiring $E_1,E_4$ gives the five units
$[1,2],3,[4,5],6,7$, whereas requiring $E_3,E_4$ gives
$1,2,[3,4,5],6,7$. Both leave five units; the overlap changes their
lengths, but not their number. The directed links force each block's
internal order, so there is no factor of two for reversing a block.

The links follow the path $1,2,\ldots,7$, so each required link reduces
the number of blocks by one. There are exactly $7-k$ distinct blocks,
each with a fixed internal order. Their $(7-k)!$ permutations give all
rows satisfying the required links. Contracting the specified links uniquely recovers the block order,
and expansion reverses the operation. If two blocks create an extra
success at their boundary, they remain separate units in this count:
only specified links are contracted. Extra successes are allowed in
this intermediate count. Every number is used once, so no indistinguishability
factor is needed.

\textbf{Require two specified successes and exclude the other four.}
Fix a two-element set $S$ of success indices. There are $5!$ permutations
satisfying its two required links. To make these the \emph{only} successes,
apply inclusion--exclusion to the four indices outside $S$:
\[
\begin{gathered}
5!-\binom41 4!+\binom42 3!-\binom43 2!+\binom44 1!\\
=120-96+36-8+1=53.
\end{gathered}
\]
\clearpage
\solution{10 (continued)}
The following table expands the inclusion--exclusion calculation for a fixed
two-element success set $S$.
\begin{center}
\renewcommand{\arraystretch}{1.15}
\begin{tabular}{@{}ccccr@{}}
\toprule
Extra links $j$ & Choices & Blocks & Orders per choice & Signed term \\
\midrule
$0$ & $1$ & $5$ & $5!=120$ & $+120$ \\
$1$ & $4$ & $4$ & $4!=24$ & $-96$ \\
$2$ & $6$ & $3$ & $3!=6$ & $+36$ \\
$3$ & $4$ & $2$ & $2!=2$ & $-8$ \\
$4$ & $1$ & $1$ & $1!=1$ & $+1$ \\
\bottomrule
\end{tabular}
\end{center}
Subtracting single unwanted events removes rows with any extra success,
but removes a row with two extra successes twice. Adding intersections
restores that excess subtraction; higher intersections continue the
same correction. The table tracks these alternating corrections.

Indeed, requiring any additional $j$ of the four links leaves
$(5-j)!$ block permutations, and there are $\binom4j$ ways to specify them.
The contraction argument also covers overlapping links, so these counts
apply to every choice of $S$.

To see why the correction works, a row with $r$ additional successes
has total weight
\[
\sum_{j=0}^{r}(-1)^j\binom rj=(1-1)^r.
\]
This is $0$ if $r>0$, while a row with no additional successes has weight
$1$. Thus exactly the desired rows survive.

There are $\binom62$ choices for $S$. Every row with exactly two successes
has one unique such set, so multiplying introduces no overcount. The answer is
\[
\binom62\cdot53=15\cdot53=\boxed{795}.
\]

\solution{11}
\textbf{Translate the geometric restriction into a condition on a string.}
A route is an ordered string containing six R's and six U's. The copies
of each step are indistinguishable, but their positions matter. After
any initial segment, $x$ is the number of R's used and $y$ the number of
U's used. Thus a route is valid precisely when every initial segment
contains at least as many R's as U's.

Ignoring this restriction gives:
\vspace{-8mm}
\[
\frac{12!}{6!\,6!}=\binom{12}{6}=924
\]
step sequences. We will subtract those that go above the diagonal.

\textbf{Reflect the prefix ending at the first violation.}
Take any forbidden route and stop the first time it reaches $y=x+1$.
Such a time exists: a single step changes $y-x$ by one, so a route
cannot go from $y\leq x$ to $y>x$ without reaching that line.
If this prefix has $a$ R's, it has $a+1$ U's.
Exchange R and U throughout this prefix and leave the rest unchanged.
Geometrically, this reflects the prefix across $y=x$:
\[
\underbrace{\text{prefix: }a\text{ R's},\ (a+1)\text{ U's}}_{
\text{ends at the first }y=x+1}
\quad\longleftrightarrow\quad
\underbrace{\text{prefix: }(a+1)\text{ R's},\ a\text{ U's}}_{
\text{ends at the first }x=y+1}.
\]
The whole transformed string has seven R's and five U's: one U has
become an R overall. There is no restriction on the rest of this string.

\textbf{Prove that the construction is reversible.}
Start with any string of seven R's and five U's. Since it ends with
$x-y=2$, it must have a first prefix ending at $x=y+1$.
Swap R and U in this prefix. The resulting string has six of each step.
Before the chosen endpoint the original prefix has $x\leq y$;
after reflection this becomes $y\leq x$. At the endpoint it becomes
$y=x+1$. Thus this is exactly the first violation of the resulting
forbidden route.

Conversely, before the first violation of a forbidden route we have
$y\leq x$, so its reflected prefix has $x\leq y$ until it first reaches
$x=y+1$. Each construction therefore identifies the same stopping point
and undoes the other. This is a one-to-one correspondence: every
forbidden route is counted once, and every seven-R, five-U string
corresponds to a forbidden route.

\textbf{Count and subtract.} Choose the five positions occupied by U
in a transformed string. There are $\binom{12}{5}=792$ choices; do not
multiply by a factorial for identical steps. Valid and forbidden routes
are disjoint and exhaust the unrestricted routes, so the answer is
\[
\binom{12}{6}-\binom{12}{5}=924-792=\boxed{132}.
\]

\clearpage
\stageheading{Summarized Solutions}
\begingroup
\setlist[itemize]{leftmargin=*,itemsep=0pt,topsep=3pt,parsep=0pt}

\solution{1}
\[
6!=\boxed{720}.
\]
\begin{itemize}
\item \(6\): total distinct shelf books.
\item \(720\): distinct possible book arrangements.
\end{itemize}

\solution{2}
\[
\binom94=\boxed{126}.
\]
\begin{itemize}
\item \(9\): available distinct club members.
\item \(4\): members on each committee.
\item \(126\): distinct possible member committees.
\end{itemize}

\solution{3}
\[
8\cdot7\cdot6=\boxed{336}.
\]
\begin{itemize}
\item \(8\): first place winner choices.
\item \(7\): second place winner choices.
\item \(6\): third place winner choices.
\item \(336\): distinct possible award assignments.
\end{itemize}

\solution{4}
\[
\frac{7!}{2!\,2!}=\boxed{1260}.
\]
\begin{itemize}
\item \(7\): total letters to arrange.
\item \(2\) (first denominator): indistinguishable copies of L.
\item \(2\) (second denominator): indistinguishable copies of O.
\item \(1260\): distinct possible letter strings.
\end{itemize}

\clearpage
\solution{5}
\[
\frac{5!}{2!}\binom62=\boxed{900}.
\]
\begin{itemize}
\item \(5\): letters remaining without L.
\item \(2\) (factorial): indistinguishable copies of O.
\item \(6\): available gaps for insertion.
\item \(2\) (binomial): separate gaps receiving L's.
\item \(900\): strings with separated L's.
\end{itemize}

\solution{6}
\[
2\cdot5!-4\cdot4!=\boxed{144}.
\]
\begin{itemize}
\item \(2\): internal AB block orders.
\item \(5\): units with AB contracted.
\item \(4\) (coefficient): combined pair orientation choices.
\item \(4\) (factorial): units with both contracted.
\item \(144\): arrangements satisfying both restrictions.
\end{itemize}

\solution{7}
\[
36+648=\boxed{684}.
\]
\begin{itemize}
\item \(36\): rows containing two girls.
\item \(648\): rows containing three girls.
\item \(684\): total valid student rows.
\end{itemize}

\clearpage
\solution{8}
\[
\binom83=\boxed{56}.
\]
\begin{itemize}
\item \(8\): available compressed candidate numbers.
\item \(3\): numbers in each subset.
\item \(56\): valid nonconsecutive number subsets.
\end{itemize}

\solution{9}
\[
6\binom83=\boxed{336}.
\]
\begin{itemize}
\item \(6\): ordered run length patterns.
\item \(8\): stars and bars positions.
\item \(3\): dividers separating four gaps.
\item \(336\): subsets satisfying pair restriction.
\end{itemize}

\solution{10}
\[
\binom62\cdot53=\boxed{795}.
\]
\begin{itemize}
\item \(6\): available consecutive success indices.
\item \(2\): required consecutive successes.
\item \(53\): permutations per success set.
\item \(795\): permutations with two successes.
\end{itemize}

\solution{11}
\[
\binom{12}{6}-\binom{12}{5}=\boxed{132}.
\]
\begin{itemize}
\item \(12\): total route step positions.
\item \(6\): rightward steps per route.
\item \(5\): upward steps after reflection.
\item \(132\): routes never above diagonal.
\end{itemize}
\endgroup

\end{document}
